\setcounter{section}{1}
\phantomsection
\section*{Abstract}

This paper examines the influence of birth order on scholarly productivity among historical European scholars. Using a new dataset of 645 brother pairs from the Repertorium Eruditorum Totius Europae (RETE), we compare the academic output of older and younger siblings within families. The empirical strategy relies on within-family comparisons, helping to control for shared family characteristics. Scholarly productivity is measured using bibliographic indicators from the Virtual International Authority File (VIAF) together with supplementary measures of academic visibility. The results indicate a modest productivity advantage for older brothers: they are more likely to have at least one recorded publication and tend to display slightly higher levels of scholarly output. However, conditional on publishing, productivity differences become substantially smaller. The results also reveal substantial historical heterogeneity. Before the sixteenth century, productivity differences are more strongly associated with sibling age gaps, whereas after the sixteenth century differences in institutional affiliation become more strongly associated with productivity divergence between brothers. Overall, the findings suggest that differences in institutional affiliation are more consistently associated with within-family productivity divergence than regional inheritance systems.

%\section*{Resumo em português}

%Este artigo analisa a influência da ordem de nascimento na produtividade académica, em termos históricos, entre académicos europeus. Utilizando um novo conjunto de dados com 645 pares de irmãos do Repertorium Eruditorum Totius Europae (RETE), comparamos a produção académica de irmãos mais velhos e mais novos no seio da mesma família. A estratégia empírica assenta em comparações intrafamiliares, permitindo controlar características familiares comuns aos irmãos. A produtividade académica é medida através de indicadores bibliográficos do Virtual International Authority File (VIAF), complementados por medidas adicionais de visibilidade académica. Os resultados apontam para uma vantagem moderada dos irmãos mais velhos: apresentam maior probabilidade de ter pelo menos uma publicação registada e tendem a revelar níveis ligeiramente superiores de produção académica. No entanto, condicionando a análise à existência de publicação, as diferenças de produtividade tornam-se substancialmente mais reduzidas. Os resultados revelam ainda uma heterogeneidade histórica significativa. Antes do século XVI, as diferenças de produtividade estão mais associadas às diferenças etárias entre irmãos; após o século XVI, as diferenças de afiliação institucional passam a estar mais fortemente relacionadas com a divergência de produtividade académica. De forma geral, os resultados sugerem que os contextos institucionais desempenharam um papel mais relevante do que os sistemas de herança na explicação das diferenças de produtividade intrafamiliares entre académicos.